\begin{abstract}
\vspace{-1em}
Inverse rendering of glossy objects from RGB imagery remains fundamentally limited by inherent ambiguity. Although NeRF-based methods achieve high-fidelity reconstruction via dense-ray sampling, their computational cost is prohibitive. Recent 3D Gaussian Splatting achieves high reconstruction efficiency but exhibits limitations under specular reflections. Multi-view inconsistencies introduce high-frequency surface noise and structural artifacts, while simplified rendering equations obscure material properties, leading to implausible relighting results. To address these issues, we propose GOGS, a novel two-stage framework based on 2D Gaussian surfels. First, we establish robust surface reconstruction through physics-based rendering with split-sum approximation, enhanced by geometric priors from foundation models. Second, we perform material decomposition by leveraging Monte Carlo importance sampling of the full rendering equation, modeling indirect illumination via differentiable 2D Gaussian ray tracing and refining high-frequency specular details through spherical mipmap-based directional encoding that captures anisotropic highlights. Extensive experiments demonstrate state-of-the-art performance in geometry reconstruction, material separation, and photorealistic relighting under novel illuminations, outperforming existing inverse rendering approaches.
\end{abstract}